\clearpage\pagestyle{otMain}\setlength{\otBlocEspacement}{2.5mm}\fontsize{11}{15}\selectfont\corrigesChapitre{11}{Comparer et mesurer des aires}
\begin{blocLivre}{cor-11-1}\begin{corrige}{1}{Comparer sans mesurer}\demarche{Compter les carreaux et assembler les moitiés.}
A : 8 carreaux. B : 7 carreaux. C : 6 carreaux et deux demi-carreaux, soit 7 carreaux. Donc $\res{B\text{ et }C}$ ont la même aire, plus petite que celle de A.
\label{fig-cor-11-1}\figureGeo[corrige]{air-1}
\end{corrige}\end{blocLivre}
\begin{blocLivre}{cor-11-2}\begin{corrige}{2}{Compter les carreaux}\demarche{Compléter en rectangle, puis retirer.}
Le rectangle de 5 sur 3 compte 15 carreaux ; il en manque 2. Aire : $15-2=\res{13\ \mathrm{cm}^2}$.
\label{fig-cor-11-2}\figureGeo[corrige]{air-2}
\end{corrige}\end{blocLivre}
\begin{blocLivre}{cor-11-3}\begin{corrige}{3}{Choisir l’unité}\demarche{Imaginer le carré unité à côté de l’objet.}
Timbre : $\res{\mathrm{cm}^2}$. Page de cahier : $\res{\mathrm{dm}^2}$, environ 6 dm$^2$. Chambre : $\res{\mathrm m^2}$, environ 10 m$^2$.
\end{corrige}\end{blocLivre}
\begin{blocLivre}{cor-11-4}\begin{corrige}{4}{Avec les formules}\demarche{Remplacer les lettres par les mesures.}
Carré : $7\times7=\res{49\ \mathrm{cm}^2}$. Rectangle : $12\times4{,}5=\res{54\ \mathrm{m}^2}$.
\end{corrige}\end{blocLivre}
\begin{blocLivre}{cor-11-5}\begin{corrige}{5}{Convertir}\demarche{Le rapport entre unités voisines est 100.}
$\res{300\ \mathrm{dm}^2}$ ; $\res{500\ \mathrm{cm}^2}$ ; $\res{4\ \mathrm{dm}^2}$ ; $\res{2{,}5\ \mathrm{m}^2}$.
\end{corrige}\end{blocLivre}
\begin{blocLivre}{cor-11-6}\begin{corrige}{6}{Conversions décimales}\demarche{Multiplier ou diviser par 100.}
$\res{70\ \mathrm{dm}^2}$ ; $\res{0{,}45\ \mathrm{dm}^2}$ ; $\res{320\ \mathrm{cm}^2}$ ; $\res{0{,}08\ \mathrm{m}^2}$.
\end{corrige}\end{blocLivre}
\begin{blocLivre}{cor-11-7}\begin{corrige}{7}{Même aire, autre périmètre}\demarche{Calculer séparément les deux grandeurs.}
Rectangle de 6 cm sur 2 cm : aire $\res{12\ \mathrm{cm}^2}$, périmètre 16 cm. Rectangle de 4 cm sur 3 cm : aire $\res{12\ \mathrm{cm}^2}$, périmètre 14 cm. Même aire, périmètres différents.
\label{fig-cor-11-7}\figureGeo[corrige]{air-7}
\end{corrige}\end{blocLivre}
\begin{blocLivre}{cor-11-8}\begin{corrige}{8}{Découper et recoller}\demarche{Transformer la figure en rectangle sans changer l’aire.}
Le triangle recollé à droite complète un rectangle de 4 carreaux sur 3, visible en vert. L’aire de la figure de départ est donc $4\times3=\res{12\ \mathrm{cm}^2}$.
\label{fig-cor-11-8}\figureGeo[corrige]{air-8}
\end{corrige}\end{blocLivre}
\begin{blocLivre}{cor-11-9}\begin{corrige}{9}{Retrouver une longueur}\demarche{Chercher le facteur manquant.}
$9\times\square=36$ : largeur $\res4$ cm. $c\times c=64$ : $c=\res8$ m, car $8\times8=64$.
\end{corrige}\end{blocLivre}
\begin{blocLivre}{cor-11-10}\begin{corrige}{10}{La pièce en L}\demarche{Découper en rectangles, ou compléter puis retirer.}
Découpage : $9\times3+4\times4=27+16=\res{43\ \mathrm{m}^2}$. Complément : $9\times7-5\times4=63-20=43\ \mathrm{m}^2$. Les deux méthodes concordent.
\label{fig-cor-11-10}\figureGeo[corrige]{air-10}
\end{corrige}\end{blocLivre}
\begin{blocLivre}{cor-11-11}\begin{corrige}{11}{Le carrelage}\demarche{Convertir l’aire dans l’unité du carreau.}
Aire : $4\times3{,}5=14\ \mathrm{m}^2$, soit $14\times100=1\,400\ \mathrm{dm}^2$. Chaque carreau couvre 1 dm$^2$ : $\res{1\,400}$ carreaux.
\end{corrige}\end{blocLivre}
\begin{blocLivre}{cor-11-12}\begin{corrige}{12}{La peinture}\demarche{Retirer l’aire de la porte, puis interpréter la division.}
Mur : $12{,}5$ m$^2$. Porte : $1{,}8$ m$^2$. À peindre : $\res{10{,}7\ \mathrm{m}^2}$. Deux pots couvrent 10 m$^2$, c’est trop peu : il faut $\res3$ pots.
\end{corrige}\end{blocLivre}
\begin{blocLivre}{cor-11-13}\begin{corrige}{13}{Le jardin partagé}\demarche{Calculer les aires, puis la proportion.}
Jardin : 400 m$^2$. Potager : $20\times8=\res{160\ \mathrm{m}^2}$. Pelouse : $400-160=\res{240\ \mathrm{m}^2}$. Proportion : $\frac{160}{400}=\frac{40}{100}$, soit $\res{40\text{~\%}}$.
\label{fig-cor-11-13}\figureGeo[corrige]{air-13}
\end{corrige}\end{blocLivre}
\begin{blocLivre}{cor-11-14}\begin{corrige}{14}{Des timbres sur une affiche}\demarche{Exprimer les deux aires dans la même unité.}
Timbre : $3\times2{,}5=7{,}5$ cm$^2$. Affiche : $6\times4=24$ dm$^2$, soit $2\,400$ cm$^2$. Nombre : $2\,400:7{,}5=\res{320}$. Contrôle : 20 timbres sur 60 cm et 16 timbres sur 40 cm, $20\times16=320$.
\end{corrige}\end{blocLivre}
\begin{blocLivre}{cor-11-15}\begin{corrige}{15}{Vrai ou faux}\demarche{Tester sur un exemple.}
Côté 3 cm : aire 9 cm$^2$. Côté 6 cm : aire 36 cm$^2$, quatre fois plus. C’est $\res{\text{faux}}$ : l’aire est multipliée par $\res4$.
\end{corrige}\end{blocLivre}
